\documentclass[11pt]{article}
\usepackage{mathblatt}
\usepackage{multicol}


\begin{document}
\pfheftstil
\fancyfoot[C]{\parbox[t]{\textwidth}{\mbfhbox\par\vspace{1.5mm}{\footnotesize\color{mbgrau}{\scriptsize 3Z2}\hfill\thepage}}}
\pfheftkopf{Grundwert \textperiodcentered{} Lösungen}{}
{\small\noindent\textbf{Typische Fehler}\par\smallskip\noindent \textbullet\ Grundwert gesucht, aber „p \% von W“ gerechnet: 25 \% von 200 € statt 200 · 4; 20 \% von 6 €.\par\noindent \textbullet\ Grundwert und Prozentwert vertauscht: 70 : 30 statt 0,3 · 70; 7,8 : 1,25 statt 1,25 : 7,8.\par\noindent \textbullet\ Falsches Ganzes: Teil zu Rest statt Teil zu Ganzem (2 : 14 statt 2 : 16; 9 zu 6 statt 9 zu 15).\par\noindent \textbullet\ „um“ und „auf“ verwechselt: „auf 70 \% herabgesetzt“ als 70 \% Rabatt gelesen.\par}\medskip
\begin{pfloesung}{}
\lz{1.}{$G = W \cdot  100 : p$}{}{}
\end{pfloesung}
\begin{pfloesung}{}
\lz{2.}{a) $28$ Kinder; b) $150$ Gäste; c) der alte Preis, $60\,€$}{}{}
\end{pfloesung}
\begin{pfloesung}{}
\lz{3.}{a) W = 9 Kinder, p = 30; b) W = 18 €, p = 15}{}{}
\end{pfloesung}
\begin{pfloesung}{}
\lz{4.}{$5$ Abschnitte bis $100\,\%$; das Ganze ist der ganze Streifen}{}{}
\end{pfloesung}
\begin{pfloesung}{}
\lz{5.}{$240$ kg}{wie oft bis hundert Prozent: $100 : 10 = 10$; mal nehmen: $10 \cdot 24 = 240$ kg}{}
\end{pfloesung}
\begin{pfloesung}{}
\lz{6.}{$800\,€$}{$1\,\% = 8\,€$; $100\,\%$}{}
\end{pfloesung}
\begin{pfloesung}{}
\lz{7.}{$225$ kg}{$81 : 36 = 2{,}25$; $2{,}25 \cdot 100$}{}
\end{pfloesung}
\begin{pfloesung}{}
\lz{8.}{800 €}{G = 200 € · 100 : 25}{}
\end{pfloesung}
\begin{pfloesung}{}
\lz{9.}{30 €}{G = 6 € · 100 : 20}{}
\end{pfloesung}
\begin{pfloesung}{}
\lz{10.}{G = 4 : $\tfrac{2}{3}$ = 6 Kugeln ($\tfrac{1}{3}$ sind 2 Kugeln)}{}{}
\end{pfloesung}
\end{document}